\scanpage{238}

this subgroup contains $\Gamma_{\mathbb Q}$, we therefore find a section of the extension $\mathcal M_{0,3}^{!}$ of $\Gamma_{\mathbb Q}$ by $\widehat\Pi_{0,3}$, \deleted{It is necessary} \inserted{which is \unsure{in|in}} fact a section of the subgroup of $\mathcal M_{0,3}^{!}$ which is the normalizer of $L_0$. We should check that this canonical section is associated with the choice of a uniformizer \unsure{at the} point $0$ of $\widehat U_{0,3\mathbb Q}$, \deleted{\unsure{and which will be determined}} \inserted{whose image in} $\unsure{\mathfrak m_0}/\mathfrak m_0^2$ should be \inserted{well} determined, where $\mathfrak m_0$ denotes the maximal ideal of the local ring at $0$ in $\widehat U_{0,3\mathbb Q}=\mathbb P_{\mathbb Q}^1$ \unsure{)} --- we should determine this element of $\mathfrak m_0/\mathfrak m_0^2$ --- We feel like taking an element which is, not only a basis over $\mathbb Q$, but even (in the situation over the integers) a basis over $\mathbb Z$ --- \deleted{\unsure{this}} \inserted{condition} which indeed determines an element of $\mathfrak m_0/\mathfrak m_0^2$ \uuline{up to} \uline{sign} --- \unsure{this would be} in fact the choice corresponding to the uniformizer $z$, or $-z$. But is it really the right one? We should come back to this!
\margin{237}

\scanpage{239}
\mspage{420}

Let us work henceforth in the subgroup $\mathcal M_{0,3}^{!\approx}$ of $\mathcal M_{0,3}^{!\sim}$, consisting of the $u$ which commute as outer actions with $\widehat\varphi$, i.e. inverse image of $\mathcal N_{0,3}^{!\approx}$, so that we have
\[
\diagram{\begin{tikzcd}[ampersand replacement=\&]
1 \arrow[r] \& \widehat\Pi_{0,3} \arrow[r] \& \mathcal M_{0,3}^{!\approx} \arrow[r] \& \underset{\substack{\cup\\\Gamma_{\mathbb Q}}}{\mathcal N_{0,3}^{!\approx}} \arrow[r] \& 1
\end{tikzcd}}
\tag{24}
\]
We have therefore found a section subgroup
\[
\mathcal M_{0,3}^{!\approx}(0)\xrightarrow{\sim}\mathcal N_{0,3}^{!\approx}
\tag{25}
\]
\deleted{$\{\unsure{\xi}\mid\xi\in\widehat\Pi_{0,3}\text{ solution of }(14)(15)\}$ \illegible{2}}
\[
\{u_\xi\mid\xi\in\widehat\Pi_{0,3}\text{ solution of }(14)\text{ and }(15)\}\subset\mathcal M_{0,3}^{!\approx}.
\]

In terms of this we can therefore write $\mathcal M_{0,3}^{!\approx}$ as a $\frac12$ direct product
\[
\mathcal M_{0,3}^{!\approx}\simeq\mathcal M_{0,3}^{!\approx}(\unsure{0})\cdot\widehat\Pi_{0,3},
\tag{26}
\]
by writing every element of $\mathcal M_{0,3}^{!\approx}$ uniquely as
\[
\begin{gathered}
\deleted{\unsure{u_{\xi,f}}}\operatorname{int}(f^{-1})\circ\unsure{u_0|u_\xi}\\
\bigl(\unsure{u_0|u_\xi}\in\mathcal M_{0,3}^{!\approx}(0)\text{ i.e. }u_0=u_\xi\text{ with }\xi\text{ solution of }(14)(15)\bigr)\\
f\in\widehat\Pi_{0,3}\text{ well determined}
\end{gathered}
\tag{27}
\]
These are precisely the \deleted{$u_{\xi,f}=u_{\alpha,\beta,f}$} \inserted{$u_{\xi,f}$ of (19), i.e.,}\margin{\illegible{1}}
\[
u_{\xi,f}=u_{\alpha,\beta,f}
\tag{28}
\]
with $\alpha,\beta$ determined in terms of $\xi$ by
\[
\alpha\rho(\alpha)^{-1}=\xi,\quad\deleted{\unsure{\varepsilon}}\beta\sigma_\infty(\beta)^{-1}=\xi \ell_1^{-\nu},
\tag{29}
\]
\margin{238}

\scanpage{240}

this is how we can interpret the \uline{normalized} \unsure{triples} $(\alpha,\beta,f)$ of corollary 1.

\uline{\unsure{Rescholium}} \inserted{(Recapitulation)} The group $\mathcal N_{0,3}^{!\approx}\simeq\mathcal M_{0,3}^{!\approx}(0)$ is in $1$-$1$ correspondence with \unsure{the set} of $\xi\in\widehat\Pi_{0,3}$ satisfying the two relations
\[
\left\{
\begin{aligned}
\xi\rho(\xi)\rho^2(\xi)&=1\\
\xi&=\bigl(\widehat\varphi(\xi)\sigma_\infty(\widehat\varphi(\xi))^{-1}\bigr)\ell_1^\nu
\end{aligned}
\right.
\qquad\begin{gathered}\text{for}\\\nu\in\widehat{\mathbb Z}\text{ suitable}\end{gathered}
\tag{30}
\]
or again (putting $\xi=\alpha\rho(\alpha)^{-1}$, $\alpha\in\unsure{\widehat{\widetilde\Pi}_{0,3}}$) with \unsure{the set} of $\alpha\in\unsure{\widehat{\widetilde\Pi}_{0,3}}$ satisfying the \inserted{second} relation (30), in which we replace $\xi$ by $\alpha\rho(\alpha)^{-1}$. The element $u_\xi$ of $\mathcal M_{0,3}^{!\approx}(0)$ defined by $\xi$ is defined by the formulas (18), i.e. by
\[
u_\xi(\sigma_\infty)=(\xi \ell_1^{-\nu})\sigma_\infty,
\qquad u_\xi(\rho)=\xi\rho.
\tag{31}
\]
Finally, if we have $u_\xi$ and $u_{\xi'}$, then
\[
u_{\xi'}u_\xi=u_{\xi''},
\tag{32}
\]
where $\xi''$ is given by the formulas on page \unsure{395'} (52) --- noting that here $\xi=g$, $\xi'=g'$, $\xi''=g''$:
\deleted{$\xi''=\alpha''\rho(\alpha'')^{-1}$ with $\alpha''=u_{\xi'}(\alpha)\cdot\alpha'$, \unsure{whence|$f''=1$}}
\[
\deleted{\xi''=\deleted{\unsure{\alpha''\rho}}u_{\xi'}(\alpha)\underbrace{\alpha'\rho(\alpha')^{-1}}_{\xi'}\rho(u_{\xi'}(\alpha))^{-1}}
\]
\margin{239}

\scanpage{241}
\mspage{421}

\deleted{\illegible{1}}
\[
\xi''=u'(\xi)\cdot\xi'.
\tag{33}
\]

Let us now set
\[
\mathcal N_{0,3}^{\approx}
=\mathcal G_3\times\mathcal N_{0,3}^{!\approx}
\subset\mathcal N_{0,3}^{\sim}
=\mathcal G_3\times\mathcal N_{0,3}^{!\sim}
\tag{34}
\]
and let $\mathcal M_{0,3}^{\approx}\subset\deleted{\unsure{\mathcal G}}\mathcal M_{0,3}^{\sim}$ be the inverse image in $\mathcal M_{0,3}^{\sim}$ of $\mathcal N_{0,3}^{\approx}\subset\mathcal N_{0,3}^{\sim}$, so that we have an extension structure
\[
1\longrightarrow\deleted{\unsure{\widehat{\mathcal G}_{0,3}^{+}}}\widehat\Pi_{0,3}
\longrightarrow\mathcal M_{0,3}^{\approx}
\unsure{\twoheadrightarrow|\longrightarrow}
\underset{\simeq\mathcal G_3\times\mathcal N_{0,3}^{!\approx}}{\mathcal N_{0,3}^{\approx}}
\longrightarrow1
\tag{35}
\]
containing the \deleted{extension}\inserted{\unsure{subextension}} (24) as the inverse image of \deleted{\unsure{the}}
\deleted{\illegible{7}}
\[
\mathcal N_{0,3}^{!\approx}
=\operatorname{Ker}(\mathcal N_{0,3}^{\approx}\longrightarrow\mathcal G_3).
\]
\deleted{\unsure{(36)}} \inserted{\unsure{We have}} \deleted{\unsure{also}} a \unsure{further} \unsure{extension structure on} $\mathcal M_{0,3}^{\approx}$
\[
1\longrightarrow\widehat{\mathcal G}_{0,3}^{+}
\longrightarrow\mathcal M_{0,3}^{\approx}
\longrightarrow\mathcal N_{0,3}^{!\approx}\longrightarrow1,
\tag{36}
\]
containing (24) as a subextension of the same quotient by a smaller group.

The subgroup $\mathcal M_{0,3}^{!\approx}(0)$, which is a section of (24), is therefore also a section of (36), hence a canonical isomorphism
\[
\mathcal M_{0,3}^{\approx}
\simeq\underbrace{\unsure{\mathcal N|\mathcal M}_{0,3}^{!\approx}(0)}_{\simeq\mathcal N_{0,3}^{!\approx}}
\cdot\widehat{\mathcal G}_{0,3}^{+}
\tag{37}
\]
($\frac12$-direct product).
\margin{240}

\scanpage{242}

Let us now consider the map
\[
\Phi:\quad
\diagram{\begin{tikzcd}[ampersand replacement=\&,column sep=large]
\mathcal M_{0,3}^{!\approx}\arrow[r]\arrow[d,phantom,sloped,"\simeq"]
\&\mathcal M_{0,3}^{\approx}\arrow[d,phantom,sloped,"\simeq"]\\
\unsure{\mathcal M}_{0,3}^{!\approx}(0)\cdot\widehat\Pi_{0,3}
\&\mathcal M_{0,3}^{!\approx}(0)\cdot\widehat{\mathcal G}_{0,3}^{+}
\end{tikzcd}}
\tag{38}
\]
which is \unsure{given|given} by
\[
\Phi(u\cdot a)=u\widehat\varphi(a)
\tag{39}
\]
($u\in\mathcal M_{0,3}^{!\approx}(0)$, $a\in\widehat\Pi_{0,3}$). It follows from the characterization of $\mathcal M_{0,3}^{!\approx}(0)$ (corollary 3) --- \deleted{\unsure{i.e.}} \inserted{more \uline{particularly}} (from the fact that the elements of $\mathcal M_{0,3}^{!\approx}(0)$ actually commute (not only as outer actions) with $\widehat\varphi$, that $\Phi$ is a group homomorphism. This is of course none other than the homomorphism $\Phi$ considered in \S\,41 (page\,316') \inserted{(up to restriction to the subgroups $\mathcal M_{0,3}^{!}=\unsure{\mathcal E}_{0,3}$ and $\mathcal M_{0,3}$)} --- in loc. cit. this homomorphism could moreover be defined in an effective way --- not only up to $\widehat\Pi_{0,3}$-conjugation or up to $\widehat{\mathcal G}_{0,3}^{+}$-conjugation --- and with a little care we should be able to verify \uline{the equality} of these homomorphisms $\Phi$ from \S\,41 and from here, in an \uline{effective} way.
\margin{241}

\scanpage{243}
\mspage{422}
\uline{Calculation \unsure{with|of}} $x=\ell_\infty$.

I am going to return to the determination of the $u_\xi\in\mathcal M_{0,3}^{!\sim}[0]=\operatorname{Norm}_{\mathcal M_{0,3}^{!\sim}}(L_0)$ which are \uline{normalized} i.e. which commute with $\widehat\varphi$ i.e. satisfy (5) with $\lambda=1$. For $x=\ell_0$ this formula, by virtue of the particular form of $u_\xi$, is tautological
\[
(\widehat\varphi u\ell_0=\widehat\varphi(\ell_0^\mu)=\varepsilon_0^\mu,
\quad u\widehat\varphi(\ell_0)=u(\varepsilon_0)=\varepsilon_0^\mu),
\]
so we \unsure{want|see} that it \unsure{should come|comes} to the same thing to impose it for $x=\ell_1$, or for $x=\ell_\infty$. For $x=\ell_1$, the condition takes the form (10) --- remarkable in that it gives a calculation \unsure{of the|of} ``delicate'' \deleted{the \unsure{variable}}\inserted{\unsure{parameter}} $\beta$ (whose existence \unsure{au priori|a priori}, in terms of $\xi$, rests on foundations which have not yet been established\ldots), in terms of the ``coarse'' parameter $\xi$ --- taking into account the fact that $i$\inserted{$\in\mathbb Z/2\mathbb Z$} itself is determined just as easily in terms of $\xi$.

We are going to give this condition a form rather different in appearance, by writing the commutation relation (5) for $x=\ell_\infty$, instead of $x=\ell_1$. Let us first note that we have the relation
\[
\tau'_\infty(\rho)=\rho^{-1}
\tag{40}
\]
For
\[
\tau'_\infty(\rho)
=\sigma_\infty(\underbrace{\tau_\infty(\rho)}_{\ell_1^{-1}\rho})
=\deleted{\underbrace{\unsure{(\sigma_\infty\ell_1^{-1}\sigma_\infty^{-1})}}_{\unsure{\ell_0^{-1}}}}\inserted{\unsure{\ell_0^{-1}}}
\underbrace{\sigma_\infty(\rho)}_{\underbrace{\sigma_\infty\rho}_{\varepsilon_0}\sigma_\infty^{-1}}
=\varepsilon_0^{-1}\sigma_\infty^{-1}=\rho^{-1},
\]
\margin{242}

\scanpage{244}

it follows that if $u\in\mathcal M_{0,3}^{!\sim}$ satisfies
\[
u(\rho)=\rho^{\unsure{\mu|j}}
\qquad \unsure{\mu|j}\in\mathbb Z/3\mathbb Z\text{ suitable}
\tag{41}
\]
then we necessarily have $j\equiv-1\text{ or }+1\ (3)$ and if \deleted{\unsure{$\alpha$ is} \illegible{1}$(j)=\alpha_3$ is the integer \illegible{1}}
\[
\alpha_3(j)\in\mathbb Z/2\mathbb Z
\quad\text{characterized by}\quad
j\equiv(-1)^{\alpha_3(j)}\ (3)
\tag{42}
\]
we will therefore have $u(\rho)=\rho$ or $u(\rho)=\rho^{-1}$ according as $\alpha_3(j)=0$ or $\alpha_3(j)=1$, and in any case we will have
\[
u(\rho)=\tau'_\infty{}^{\alpha_3}(\rho).
\tag{43}
\]

This being so, let us put $x=\ell_\infty$ in (5):
\[
\widehat\varphi u(\ell_\infty)
=\widehat\varphi(\operatorname{int}(f_\infty^{-1})\ell_\infty^\mu)
=\operatorname{int}(\widehat\varphi(f_\infty^{-1}))\rho^{-\mu}
\]
(since $\widehat\varphi(\ell_\infty)=\rho^{-1}$), and
\[
u(\widehat\varphi(\ell_\infty))
=u(\rho^{-1})=u(\rho)^{-1}
=\deleted{\unsure{(\xi\rho)^{-1}}=\rho^{-1}\xi^{-1}
=\unsure{(\rho^{-1}\xi^{-1}\rho)\rho^{-1}}}
\inserted{\operatorname{int}(\alpha)(\rho^{-1})}
\]
so formula (5) for $\lambda=1$ becomes
\[
\operatorname{int}(\widehat\varphi(f_\infty)\alpha)(\rho^{-1})
=\rho^{-\mu}\qquad\text{i.e.}
\]
\[
\operatorname{int}(\widehat\varphi(f_\infty)\alpha)\rho=\rho^\mu
\]
and by virtue of (43) \deleted{\unsure{this is equivalent} $\unsure{\rho^\mu=\tau'_\infty{}^{\alpha_3}(\rho)}$}\inserted{\unsure{we have}} \deleted{\unsure{that}}
\[
\rho^\mu=\tau'_\infty{}^{\alpha_3=\alpha_3(\mu)}(\rho)
\overset{\unsure{\mathrm{df}|\mathrm{déf}}}{=}
\operatorname{int}(\tau'_\infty{}^{\alpha_3})(\rho)
\]
so the condition becomes
\[
\operatorname{int}(\tau'_\infty{}^{\alpha_3}\widehat\varphi(f_\infty)\alpha)\rho=\rho
\tag{\unsure{44}}
\]
i.e. that the element $\tau'_\infty{}^{\alpha_3}\widehat\varphi(f_\infty)\alpha$ of $\widehat{\mathcal G}_{0,3}=\widehat{\mathcal G}_{0,3}^{+}\cdot\{1,\tau_\infty\}$ commutes with $\rho$.
\margin{243}